% Generated draft: requires mathematical and source review.
\documentclass[11pt]{article}
\usepackage[margin=25mm]{geometry}
\usepackage{fontspec}
\setmainfont{Latin Modern Roman}
\usepackage{amsmath,amssymb,mathtools,unicode-math}
\setmathfont{Latin Modern Math}
\usepackage{enumitem,needspace,xurl,hyperref}
\setlength{\parindent}{0pt}
\setlength{\parskip}{5pt}
\setlength{\emergencystretch}{3em}
\newcommand{\N}{\mathbb{N}}
\newcommand{\Z}{\mathbb{Z}}
\newcommand{\Q}{\mathbb{Q}}
\newcommand{\R}{\mathbb{R}}
\newcommand{\C}{\mathbb{C}}
\newcommand{\F}{\mathbb{F}}
\newcommand{\A}{\mathbb{A}}
\newcommand{\PP}{\mathbb{P}}
\newcommand{\E}{\mathbb{E}}
\newcommand{\eps}{\varepsilon}
\newcommand{\dd}{\,\mathrm{d}}
\begin{document}
\section*{UnsolvedMath ID: 30006053}
\textbf{Source problem code:} catalogue entry
\subsection{Definitions and mathematical statement}
The source asks about realization of fundamental triples and classification of Poincar\'e-duality four-dimensional complexes. It does not supply definitions of a fundamental triple or the symmetric $2$-type. The following gives a conditional formulation using the usual interpretation of a fundamental triple; the unresolved conventions are stated explicitly.

Let $T$ be a connected homotopy $2$-type, meaning a connected space with $\pi_i(T)=0$ for $i>2$, and let $w\colon\pi_1(T)\to\{\pm1\}$ be an orientation character. Write $\mathbb Z^w$ for the corresponding local system of infinite cyclic groups. A Poincar\'e-duality $4$-complex $X$ has an orientation character $w_X$ and a fundamental class $[X]\in H_4(X;\mathbb Z^{w_X})$ such that cap product induces isomorphisms
\[
-\frown[X]\colon H^k(X;M)\xrightarrow{\cong}H_{4-k}(X;\mathbb Z^{w_X}\otimes_{\mathbb Z}M)
\]
for every local coefficient system $M$ of abelian groups and every integer $k$. Any additional finiteness requirements on $X$ must be fixed; they are not specified in the supplied excerpt.

Under this interpretation, a realization of $(T,w,b)$ consists of such an $X$ and a map $c\colon X\to T$ inducing isomorphisms on $\pi_1$ and $\pi_2$, with $w_X=c^*w$ and $c_*[X]=b$. This requires
\[
b\in H_4(T;\mathbb Z^w).
\]
Define the realization set
\[
\mathcal R_{\mathrm{PD}}(T,w)=\bigl\{c_*[X]\in H_4(T;\mathbb Z^w): (X,c)\text{ satisfies these realization conditions}\bigr\}.
\]
Define $\mathcal R_{\mathrm{man}}(T,w)$ by requiring additionally that $X$ be a closed connected $4$-manifold. The category of manifolds, such as topological or smooth, is not specified by the source and must be fixed before this second set is unambiguous.

The realization questions are to determine, for every admissible pair $(T,w)$, the sets $\mathcal R_{\mathrm{PD}}(T,w)$ and $\mathcal R_{\mathrm{man}}(T,w)$, and consequently to characterize those triples $(T,w,b)$ that admit each kind of realization. Given $b\in\mathcal R_{\mathrm{PD}}(T,w)$, determine all other classes $b'$ for which $(T,w,b')$ is realized.

There are two unresolved textual qualifications. First, the source explicitly places $b'$ in $H_4(T;\mathbb Z)$, whereas the orientation-character interpretation above requires $H_4(T;\mathbb Z^w)$. These agree when $w$ is trivial; the supplied record does not state that restriction or define an alternative coefficient convention. Second, the original question quantifies over $b'$ but repeats $(T,w,b)$ as its target. The preceding formulation with $(T,w,b')$ follows the supplied cleaned statement and remains an interpretation requiring review.

The classification question asks what additional data must be attached to the symmetric $2$-type of a Poincar\'e-duality $4$-complex to determine its homotopy type in general, and whether these data can be expressed as a quadratic enhancement, called a ``quadratic $2$-type.'' The supplied excerpt defines neither the symmetric $2$-type nor the proposed quadratic enhancement or their equivalence relations. Thus this last question specifies a research objective, but does not yet specify a complete invariant or a precise classification assertion.
\subsection{Short English statement}
Determine which fundamental triples arise from Poincar\'e-duality $4$-complexes and which arise from $4$-manifolds. For fixed underlying data $(T,w)$ admitting a realization, determine all realizable fundamental classes. Also determine whether a quadratic enhancement of the symmetric $2$-type can recover the homotopy type. The coefficient system, finiteness conventions, manifold category, and classification invariants require clarification.
\subsection{Sources}
\begin{itemize}
\item Local registry: \texttt{data/unsolvedmath/problems.json}, permanent UnsolvedMath ID $30006053$, source problem code \texttt{catalogue entry}, title ``Realization and Classification of Poincar\'e Duality Four-Complexes.'' The supplied record lists proposed year $2024$.
\item Supplied source citation: \emph{MATRIX-MFO Tandem Workshop: Invariants in Low-Dimensional Topology: Combinatorics, Geometry, and Computation} ($2025$), Oberwolfach Reports, pp.\ 2511--2550, DOI \url{https://doi.org/10.4171/owr/2024/43}. The supplied literature list labels this ``Oberwolfach Report 43 (2024).'' The differing year labels are retained without reconciliation.
\item The original excerpt attributes the contribution ``Poincar\'e duality $4$-complexes'' to Beatrice Bleile, Diarmuid Crowley, Daniel Galvin, Jayden Hammet, Naageswaran Manikandan, Mark Pencovitch, and Arunima Ray. The registry supplies no separate proposer attribution.
\item Supplied literature reference: Baues and Bleile, \emph{Poincare duality complexes in dimension four}, \url{https://msp.org/agt/2008/8-4/agt-v8-n4-p18-s.pdf}.
\item Literature assessment dated exactly \texttt{2026-08-22}: the supplied curation labels the problem open and states that the report poses general realization and symmetric-$2$-type refinement questions, which the cited realization theory does not settle in full. This is supplied dated curation, not an independently verified literature assessment.
\item The registry labels the statement \texttt{corrected\_verified} and reports that the official report was checked and its notation repaired. These are supplied provenance claims; the present statement-only draft has not independently checked the report or the literature.
\end{itemize}
\end{document}
